\section{\texorpdfstring{局部 CIM 服务单元的 $\rho$、$\tau$、$\RI^*$ 估算方法}{局部 CIM 服务单元的 ρ、τ、RI* 估算方法}}

\subsection{从逻辑服务字节与服务间隔出发}

一个完整 streaming 服务单元对应 $B_S$ Byte 输入，完成规定求值与输出；一个完整 resident 更新服务单元对应 $B_R$ Byte，更新后已建立可计算的逻辑状态。它们分别以 $\Delta_S$、$\Delta_R$ 秒的服务间隔完成，则
\begin{equation}
\boxed{\rho\approx\frac{B_S}{\Delta_S},\qquad
\tau\approx\frac{B_R}{\Delta_R},\qquad
\RI^*=\frac{\rho}{\tau}\approx
\frac{B_S}{B_R}\frac{\Delta_R}{\Delta_S}.}
\label{eq:service}
\end{equation}
$\rho,\tau$ 的单位均为 Byte/s，$\RI^*$ 为同量纲服务字节的比值，数值无量纲。$B_S,B_R$ 是硬件事务粒度；Task 0 的 $Q_S,Q_R$ 则是在 workload 窗口内累计的逻辑输入与逻辑写入量，包含实际重放和重载。内部位展开、ADC 批次和编程重试只增加相应服务占用。

本参考 $B_S=n_{\rm in}b_S=128$ Byte。若一次局部更新真正完成 16 个 INT8 权重，$B_R=16$ Byte，式~\eqref{eq:service} 的字节因子为 8；若把同样的连续更新聚合为整份 16384 Byte 状态，则 $\Delta_R$ 也应包含完成整份状态的全部时间。以均匀、无额外固定开销的更新为例，分子和分母同比聚合不改变 $\tau$。因此 ridge 一般不等于“写延迟/读延迟”。相同单元独立复制 $k$ 份、两路需求和可用能力均按相同口径汇总时，$\rho$ 与 $\tau$ 同乘 $k$，ridge 可抵消这一因子；ADC 复用、写入并行宽度、不同位串行次数及专用电压准备通常只影响其中一路。

对窗口内 $Q_R>0$，接回已审阅模型：
\[
\RI=Q_S/Q_R,\qquad P=Q_S/T\leq\min(\rho,\tau\RI).
\]
$Q_R=0$ 使用静态分支 $\RI=\infty$、$P\leq\rho$。用本节的估计能力代入，得到的是对应参考情景的估计 Roofline；区间表示配置和参数的不确定性，不是统计置信区间，也不是全部实际实现的严格包络。两路可分别估计，即使共享阵列而互斥，原上界仍成立；能否同时达到其独立能力由实际调度决定。

\subsection{Streaming：把完整精度落实为物理步骤}

\textbf{ACIM 主参考。} 将 INT8 输入按 8 个 bit 展开，8 个二进制权重位平面同时求值。对每个输入 bit，每平面 128 个输出需由 16 个 ADC 分 8 批读出。因此一次向量有 64 次物理激励/求值、64 批转换，标量 ADC 结果总数为
\[
8\ \text{输入 bit}\times128\ \text{输出}\times8\ \text{权重位平面}
=8192=64\times128.
\]
这 8192 个码完成同一次输入服务，仍只计 $B_S=128$ Byte。每批服务 16 个输出，每个输出的八个权重位平面结果按正负位权合并，再按当前输入位权累加；所有 128 个输出就绪才完成服务。

令 $t_{m,A}$ 是本次物理步骤中介质读侧的额外必要占用，包括读激励后阵列 RC/积分/响应、公共接口以外的专用驱动及必要读后恢复，内部顺序按介质组织。主参考各阶段顺序执行，服务间隔直接为
\begin{equation}
\boxed{\Delta_{S,A}=64(T_I+t_{m,A}+T_A+2T_D)+2T_D .}
\label{eq:acim}
\end{equation}
各时间用相同单位后求和，代入式~\eqref{eq:service} 前统一转成秒。若各步骤物理时间不同，用步骤求和替代 $64t_{m,A}$。多个互补读必须分别计入 $t_{m,A}$ 或步数；若能一次差分 sensing，则按那一次实际服务计。八个位平面的并行已经用在步数中，吞吐不再额外乘八。

当某介质只能同时激活较少行、轮询权重平面、或使用更少 ADC，需显式改变行组数、权重组数和转换批次。广义计数可以按“输入分片数 $\times$ 行组数 $\times$ 权重组数 $\times$ 每组输出批数”组织，但各因子必须对应实际独立循环；同一个结构约束不同时在 ADC 数量和额外折扣中计两次。多位模拟输入或多级权重模式则同时改变可分辨部分和、误差条件和数字重构，按另一声明配置重算。

\Needspace{6\baselineskip}
\textbf{DCIM 主参考。} 公共实现一次处理 32 个输入项和 16 个输出，输入逐 bit、权重八位并行。每个数字计算步骤把当前输入位与权重相乘、归约 32 项，再将带输入位权的贡献加到该组 24-bit 输出累加器。4 个行组覆盖 128 项，8 个输出组覆盖 128 个结果，所以
\[
N_D=\underbrace{8}_{\text{输入 bit}}
\underbrace{\left\lceil128/32\right\rceil}_{\text{行组}}
\underbrace{\left\lceil128/16\right\rceil}_{\text{输出组}}=256.
\]
令 $t_{m,D}$ 为每组数字求值前、公共数字阶段以外的物理读侧时间及必要恢复，则
\begin{equation}
\boxed{\Delta_{S,D}=256(t_{m,D}+T_D)+2T_D .}
\label{eq:dcim}
\end{equation}
输入组选择和部分和寄存已在 $T_D$ 中，输入捕获和最终提交在两拍边界开销中。主参考完整精度计数为 256 步，不把每一拍当作服务了 128 Byte 新输入。若实现的权重也按位串行，应按实际独立循环增加权重步骤；若采用 2-bit/4-bit 输入、整字并行乘法或更多输出通道，应重算步数和所需时钟，而不是保持原拍数再额外乘并行收益。

文献宏若已有完整计算周期 $T_{\rm cycle,full}$，并覆盖对应的物理响应、感测、逻辑和恢复，则用它替代式~\eqref{eq:dcim} 中相应的整体步骤。所引用的逻辑维度、输入位宽和输出粒度还要匹配当前分组；例如 \sid{SDCIM-01} 的 64 拍服务的是其 $128\times16$ VMM，不能直接给当前 $128\times128$ 矩阵也填 64 拍。类似地，ACIM 的整体宏周期若已包含输入、ADC 和重构，按其实际模式换算完整服务后使用，原有阶段不再追加。采用整体证据前，还须核对其资源、精度和时隙与所选共同情景兼容；不兼容的实测宏用于独立交叉检查或显式例外，不据此静默覆盖共同节拍。

\textbf{主调度与流水线对照。} 式~\eqref{eq:acim}--\eqref{eq:dcim} 的 sum 来自已选择的顺序组织，故无跨向量重叠时服务间隔等于完整串行服务时长。若局部缓冲允许输入/阵列阶段、ADC 阶段、数字阶段流水，令每个逻辑向量对共享资源 $r$ 的必要占用为 $D_r$，则稳态间隔满足 $\Delta_S\geq\max_rD_r$。只有资源互相独立、缓冲足够且依赖允许时，才能将该值作为可实现的间隔估计；同一资源上多个阶段的占用先相加。转换延迟不等于 ADC initiation interval，流水线报告优先使用后者。有限向量数还包括流水线填充、排空。v0.1 数值采用串行公式；流水化若作为后续敏感性对照，另存成对情景。

\subsection{Resident：按完整更新序列计算}

\textbf{A. 可直接更新型。} 在声明模式下，一组权重经一次完整更新周期后可计算，则
\begin{equation}
\tau\approx\frac{B_{\rm update}}{T_{\rm update,full}}.
\label{eq:direct}
\end{equation}
先由映射求完成的逻辑权重数 $L_{\rm done}$，再取 $B_{\rm update}=L_{\rm done}b_R$。若每逻辑权重需要 $c$ 个 cell（已含二进制位切片、互补与编码），一次并行 $p$ 个 cell 完成完整目标状态，且分组对齐，可取 $L_{\rm done}=\lfloor p/c\rfloor$。若这些 cell 只是一个权重的部分位，则先累计完成所有必要位再确认逻辑 payload；不要按已触及的权重数记作全部完成。

主参考无冗余二进制映射中，一个权重需 8 个 cell。128 路驱动若能同时完成 128 个相应 cell，便完成 16 Byte；若高压/电流预算只支持较少并行 cell，增加局部写步骤或减小每事务 $B_R$。数据传输、命令、编程脉冲、建立及恢复共同构成 $T_{\rm update,full}$。对于一次数据批、单次完整物理更新的简化模式，可写 $T_{\rm update,full}=2T_D+t_{\rm physical,full}$；原文的完整更新周期若已包括公共准备与完成，直接使用该周期。

\textbf{B. 分步编程与独立擦除型。} 将实际状态转换组织为预处理、所需编程步骤、读回校验、恢复及适用擦除。例如
\[
T_{\rm update,full}=T_{\rm front}
+\sum_j\bigl(T_{{\rm drive+program},j}+T_{{\rm verify},j}
+T_{{\rm recover},j}\bigr)+T_{\rm erase,amort}.
\]
这些是粗粒度的操作块；相互重叠时按资源服务能力组织。SET 和 RESET 时间不同，只说明状态转换有方向性；每次事务执行哪一步、是否先全域 RESET，再选定 SET，取决于更新策略。初次形成若为部署初始化只在相应窗口计入；每次都需要的预处理则进入每次更新。

对于单个局部编程域的页/块模式，持续整块循环重写且每次擦除后使用 $N_{p/e}$ 个有效页服务时，可以采用
\begin{equation}
\tau\approx\frac{B_{\rm page,logical}}
{T_{\rm front}+T_{\rm program,full}+T_{\rm erase}/N_{p/e}}.
\label{eq:page}
\end{equation}
$T_{\rm front}$ 是原文 program cycle 外部仍需的命令、数据装入及完成接口时间。$B_{\rm page,logical}$ 是本逻辑映射在此页事务中完成的\emph{有效 resident Byte}；物理页未使用的字节和为旧数据保留的内容进入实际成本，不扩充逻辑写入分子。$N_{p/e}$ 指每次擦除真正支撑的页服务数，若只更新部分页，按其情景重新摊销。该式中的页容量不赋予整 die 的多 plane 并行；如果页天然覆盖当前矩阵以外的内容，本单元承担的页操作及有效字节必须同时说明。保持原生 page/block 粒度即可，无需设想器件支持逐字编程。

在已有预擦除空间上 append，可将该窗口的擦除项设为零，并明确预擦除初始状态；持续循环重写则计入擦除摊销。两种模式分别报告，避免混用 append 的时间与循环重写的容量供给条件。若映射跨多个物理页才能完成一组逻辑权重，用这些页的完整服务总时间与对应逻辑字节配对。

\textbf{C. 闭环校验。} 若原文报告从发起编程至目标状态/误差终点完成的 program cycle，且已含脉冲、verify/retry，直接以该完整周期作为操作块。若只报告单脉冲 $t_p$，则需写出脉冲数与每次读回/比较时间；例如固定 $K$ 次闭环且每次后都校验的简化情景为 $K(t_p+t_v+t_r)$，另加一次性准备。$K$ 由原文分布、终止规则或明确情景给出，末次是否需要校验按流程处理。分子仍为只完成一次的 $B_R$。读回用 binary sensing 还是多位 ADC、是否共享 streaming ADC、并行比较宽度多少，均应和同一局部资源配置衔接。

\subsection{维护、逻辑精度与局部边界的放置规则}

refresh 和破坏性读后的 restore 属于内部维护：与每次求值绑定的 restore 放入该步的物理服务时间；周期刷新占用同一资源时，可将已知刷新占用比 $f$ 作用为该资源可用率 $1-f$，或直接在排程中加入刷新块。两种方法选其一。若读写受到不同资源占用，则分别作用于两路；刷新数据不作为新 workload resident payload。endurance 和寿命单列为部署约束，不混入瞬时 $\tau$。

逻辑 INT8 不规定一个 cell 存八位。每个 cell 的状态位数、互补数量、串行步骤和可同时完成的逻辑权重数，一起决定 $B_R$ 与 $\Delta_R$。位切片并行已经进入逻辑并行度时，不再重复扩充 payload；重复物理步骤已经进入时间时，不再对 Byte 额外折扣。页/块容量和 3D 层数描述组织与容量，读写并行数由选通、驱动、感测及编程资源决定。

两路能力必须引用同一份局部状态实现、器件模式和外围资源预算。可以分别测算单独提供 streaming 或更新时的服务能力；同时运行争用同一 ADC、写驱动或阵列时，记录互斥/共享条件。若更新过程需要不同电压、重新校准或较长恢复，这些代价放在相应事务末端，使“完成更新”确实意味着能够继续计算。

\subsection{情景范围、两个示意算例与后续接口}

范围生成遵循固定次序：先固定逻辑配置、映射与服务模式；给共同外围及介质时间各自定义范围；枚举物理兼容的组合；在每个组合内导出 $(\rho,\tau,\RI^*)$；保存未取整值，再取适合量级分析的展示范围，并指出主要驱动参数。共同外围采用短时隙 $(T_I,T_A,T_D)=(2,10,2)$ ns、参考 $(5,20,5)$ ns、长时隙 $(10,50,10)$ ns 三个共享情景。它们是工程预算，未被命名为已验证的 PVT 角。各介质采用相同规则，物理时间变化不必与 CMOS 时隙同向；若存在独立变化且模式兼容，应加入相应组合。

ridge 优先取各成对情景内比值的最小、最大值；仅用独立吞吐区间交叉除法得到的
\[
\left[\rho_{\min}/\tau_{\max},\ \rho_{\max}/\tau_{\min}\right]
\]
标为保守外包络，区别于实际计算的成对范围。仅当共因子以相同乘法形式作用于两路时抵消；例如 $\Delta_S=C_S T_D$、$\Delta_R=C_R T_D$ 可消去 $T_D$，而 $(aT_D+t_{m,S})/(bT_D+t_{m,R})$ 通常保留其影响。

以下\textbf{两个算例均为不绑定介质的算术示意}，其物理时间为人为选定的输入，仅检验服务粒度、单位与范围传播。三个配对场景是演示路径，不是允许参数空间的穷尽。

\textbf{示意一：ACIM 与一次直接更新。} 采用式~\eqref{eq:acim}，三个情景的 $t_{m,A}$ 依次为 10、25、40 ns；一次完整物理更新为 20、50、100 ns，128 个二进制 cell 同时完成 16 个 INT8 权重，即 $B_R=16$ Byte。总更新时间为物理更新加 $2T_D$。参考情景得到
\[
\Delta_S=64(5+25+20+2\times5)+2\times5=3850\ns,
\quad \Delta_R=50+2\times5=60\ns,
\]
\[
\rho\approx0.0332\GBs,\quad\tau\approx0.267\GBs,\quad
\RI^*=\frac{128}{16}\frac{60}{3850}\approx0.125.
\]
这里省略字节因子会使 ridge 错八倍。ADC 复用和八位输入展开已在 64 步中，payload 不再乘 64。

\textbf{示意二：DCIM 与页更新。} 采用式~\eqref{eq:dcim}，三个 $t_{m,D}$ 为 3、10、30 ns。有效逻辑页为 256 Byte，映射恰好占 2048 物理 bit；128-bit 接口需 16 个数据拍，命令/完成另占两拍。设完整编程周期依次为 5、10、20\,$\mu$s，擦除周期为 100、200、400\,$\mu$s，每擦除支撑 16 页服务。此处完整编程含其内部 verify，时间不再乘校验次数。参考情景为
\[
\Delta_S=256(10+5)+2\times5=3850\ns,
\]
\[
\Delta_R=(16+2)\times5\ns+10\us+200\us/16
=22.59\us.
\]
故 $\rho\approx0.0332\GBs$、$\tau\approx0.0113\GBs$，ridge 约 2.93；服务字节因子在此为 $128/256=1/2$。页字节、页数和擦除时间的单位分别参与一次换算。

% Generated by scripts/check_shared.py --emit; illustrative, no medium binding.
\begin{table}[htbp]\centering\small
\begin{tabular}{@{}llrrrrr@{}}\toprule
示意 & 情景 & $\Delta_S$ (ns) & $\Delta_R$ (ns) & $\rho$ (GB/s) & $\tau$ (GB/s) & $\RI^*$\\\midrule
1 & 短时隙 & 1668 & 24 & 0.07674 & 0.6667 & 0.1151\\
1 & 参考 & 3850 & 60 & 0.03325 & 0.2667 & 0.1247\\
1 & 长时隙 & 7700 & 120 & 0.01662 & 0.1333 & 0.1247\\
\midrule
2 & 短时隙 & 1284 & 11286 & 0.09969 & 0.02268 & 4.395\\
2 & 参考 & 3850 & 22590 & 0.03325 & 0.01133 & 2.934\\
2 & 长时隙 & 10260 & 45180 & 0.01248 & 0.005666 & 2.202\\
\bottomrule\end{tabular}
\caption{两个算例的成对情景。时间列保留算术结果用于复算，非器件测量精度；GB 为 $10^9$ Byte。}
\label{tab:examples}\end{table}
示意 1 的展示范围为 $\rho$=0.016--0.077 GB/s、$\tau$=0.13--0.67 GB/s，成对 ridge 为 0.11--0.13；独立端点外包络为 0.024--0.58。
示意 2 的展示范围为 $\rho$=0.012--0.1 GB/s、$\tau$=0.0056--0.023 GB/s，成对 ridge 为 2.2--4.4；独立端点外包络为 0.55--18。


保留数据中的所有计算位数便于复算；示意范围按两位有效数字向外取整，正文外围整向量时间按一位有效数字向外取整。若只有少数情景，表内原计算点与取整范围同时呈现，避免用显示精度暗示物理准确性。示意一主要由 ADC 批次和写周期影响，示意二由数字步骤及 program/erase 摊销影响；改变普通外围的同一参数后，两路和 ridge 必须一起重算。

后续每类介质只需提交以下最小输入：所选器件/状态与逻辑映射；读侧响应及终点条件；直接更新或预处理、编程、擦除、校验、恢复的完整步骤与时间；实际局部读写粒度和并行约束；page/block 的有效映射；专用驱动和维护；引用的共享版本与例外；兼容的成对情景、导出范围及主导因素。来源记录保留原值、单位、模式、PDF 页码和图表。共享参数由 \texttt{data/shared\_parameters.json} 的 v0.1 入口提供，审阅后冻结；例外单列，避免覆盖共同条件。至此接口已具备，本轮不建立各介质结果章节。
